%%%%%%%%%%%%%%%%%%%%%%preface.tex%%%%%%%%%%%%%%%%%%%%%%%%%%%%%%%%%%%%%%%%%
% sample preface
%
% Use this file as a template for your own input.
%
%%%%%%%%%%%%%%%%%%%%%%%% Springer %%%%%%%%%%%%%%%%%%%%%%%%%%

\preface

The primary goal of this textbook is not to serve as a standalone Calculus resource but rather as a companion guide featuring a wide array of standard examples worked out in SageMath. The book is designed to introduce students to Scientific Computing while simultaneously enhancing their understanding of Calculus and expanding their mathematical horizons.
 
Key topics are presented with a dual emphasis: first through traditional analytical methods, and then through SageMath-based computations for validation and deeper insights. This approach emphasizes conceptual understanding while reducing reliance on tedious, longhand calculations, empowering students with modern computational tools.
 
The book follows a hands-on, script-based approach with SageMath, renowned for its versatility and extensive use in scientific and mathematical applications. The included SageMath commands are structured as templates, allowing students to adapt them easily by substituting parameters relevant to specific problems.
 
The textbook originates from my 15 years of experience teaching Calculus to diverse undergraduate audiences, including students in mathematics, sciences, economics, and biology. However, its utility extends beyond these groups to include more advanced students and professionals seeking to explore scientific computing.
 
Importantly, the text assumes no prior experience with Calculus or computer programming, making it accessible to a broad audience. The integrated analytical and computational framework provides a modern perspective on the study of Calculus and its applications.
 
 

\vspace{\baselineskip}
\begin{flushright}\noindent
Place(s),\hfill {\it Firstname  Surname}\\
month year\hfill {\it Firstname  Surname}\\
\end{flushright}


